\section{Extreme value theory --- reading the tail without enough data}

\subsection{The problem}

Our conformal band widens by a tail quantile of past errors.
The 90th percentile of a 90-point sample rests on about 9 points.
Nine points is noise. One odd week moves the estimate a lot.
Extreme values are exactly where data is scarcest.
Extreme value theory (EVT) is statistics built for this corner.

\subsection{Peaks over threshold (POT)}

Do not model the whole error distribution. Only the tail.

\begin{enumerate}
  \item Pick a threshold $u$, e.g.\ the 80th percentile of scores.
  \item Keep only the \emph{exceedances}: values above $u$, minus $u$.
  \item Fit a generalized Pareto distribution (GPD) to them.
  \item Read any extreme quantile from the fitted curve.
\end{enumerate}

The GPD has two parameters. Scale $\sigma$ sets the tail's width.
Shape $\xi$ sets how heavy it is. $\xi > 0$ means a heavy tail:
spikes far beyond anything seen are still possible. Theory
(Pickands, 1975) says exceedances converge to a GPD for almost any
parent distribution. That is why the trick is general.

\subsection{Worked example}

Say 100 daily scores, threshold at the 80th percentile: 20 exceedances.
The empirical 99th percentile would rest on 1 point --- useless.
The GPD fitted on all 20 exceedances gives a smooth curve.
Read the 99th percentile from the curve. All 20 points support it,
not one. That is the whole gain: borrow strength from the mid-tail
to estimate the far tail.

\subsection{The guards we used}

A fitted tail can explode. So we added three guards:
at least 30 exceedances, else no fit.
Shape $\xi < 1$, else the fitted quantiles blow up.
Any failure $\to$ fall back to the plain empirical quantile.

\subsection{The punchline: it did not help here}

Tested 2026-07-22 against a pre-declared gate. Rejected.
Spike coverage moved +0.57 points (LEAR) and +0.0 (LGBM).
Why: our calibration window holds 2{,}160 hourly scores, not 90.
The empirical 90th percentile was already stable.
The real issue is that spike errors are \emph{conditional}:
the model does not see the spike coming on that specific afternoon.
No unconditional band width fixes that. EVT solves data scarcity.
Our problem was missing information, not scarce data.

\subsection{Interview line}

``I know when EVT helps: scarce tail data, unconditional errors.
We had neither, so the GPD tail failed its gate --- and that failure
pointed us to a conditional spike classifier instead.''
